% Compile with XeLaTeX (not pdfLaTeX). Save this file as UTF-8.
% On Overleaf, select XeLaTeX as the compiler.
\documentclass[11pt,a4paper]{article}
\usepackage[margin=1in]{geometry}
\usepackage[UTF8,scheme=plain,fontset=fandol]{ctex}
\usepackage{amsmath,amssymb}
\usepackage[hidelinks]{hyperref}

\setlength{\parindent}{0pt}
\setlength{\parskip}{5pt}

\newcommand{\eps}{\varepsilon}

\begin{document}

\begin{center}
    {\Large\textbf{CSC 6011 - Theory of Computation}}\\[4pt]
    {\large Fall 2026 \quad Homework 1}\\[4pt]
    Lectures 1--4 \quad (100 points)
\end{center}

\textbf{Instructor:} Tao Lin \hfill \textbf{TA:} Yefan Gao

% Replace the underline with your name. Chinese names are supported.
\textbf{Student name (姓名):} \underline{\hspace{7cm}}

\textbf{Due:} Wednesday, Sep 30, 2026, 11:59 pm (UTC+8).

\textbf{Submission:} Upload to Blackboard: \underline{\url{https://bb.cuhk.edu.cn/}}

\textbf{Instructions:}
\begin{itemize}
\item Replace each solution placeholder with your solution.
\item Before submission, please compile the LaTeX file to PDF. Submit the PDF.
\item You can use results established in lectures without reproving them.
\item You can draw a diagram in any way and insert it here as a figure. 
\item Give clear explanations for your constructions and proofs.
\item AI usage is allowed, but please try the problems yourself first. Write the solutions yourself. Direct copying and submission of AI output is not allowed.
\end{itemize} 

\section*{Problem 1: DFAs and regular expressions (25 points)}

Let $\Sigma=\{0,1\}$.
\begin{enumerate}
    \item[(a)] \textbf{(8 points)}
    Construct a DFA \textbf{with exactly two states} to recognize
    \[
        A=\{w\in\Sigma^* \mid w \text{ contains an even number of } 1 \text{s}\}.
    \]
    Zero is even.
    Give both a diagram for the DFA \textbf{and} the transition table.  Identify the start and accepting states, and explain what each state represents.

    \item[(b)] \textbf{(10 points)}
    Construct a DFA to recognize
    \[
        B=\{w\in\Sigma^* \mid w \text{ has equally many occurrences of } 01 \text{ and }10\}.
    \]
    Occurrences of 01 and 10 can overlap. For example, $010$ has one occurrence of each pattern.
    The language includes $\eps$, $0$, and $1$.
    Give a diagram \textbf{or} a transition table. Indicate the start and accept states. Explain why your DFA recognizes exactly $B$.

    \item[(c)] \textbf{(7 points)}
    Give regular expressions for $A$ (3 points) and $B$ (4 points). Handle strings of length zero and one correctly.
\end{enumerate}

\textbf{Solution.} \textit{Write your solution here.}

\vspace{5em}

\section*{Problem 2: Closure under deleting one symbol (15 points)}

For a language $L\subseteq\Sigma^*$, define
\[
    \mathrm{DROP\mbox{-}OUT}(L) = \{xz \mid xaz\in L,\ x,z\in\Sigma^*,\ a\in\Sigma\}.
\]
Thus, a string belongs to $\mathrm{DROP\mbox{-}OUT}(L)$ if it can be obtained from a string in $L$ by deleting \emph{exactly one} symbol.

Prove that if $L$ is regular, then $\mathrm{DROP\mbox{-}OUT}(L)$ is regular.

\textbf{Solution.} \textit{Write your solution here.}

\vspace{5em}

\section*{Problem 3: A nonregular language (20 points)}

Define language
\[
    C=\{0^i1^j \mid i,j\geq 0,\ i\neq j\}.
\]
\begin{enumerate}
    \item[(a)] \textbf{(15 points)}
    Prove that $C$ is not regular by using the pumping lemma.

    \item[(b)] \textbf{(5 points)}
    Give a second proof using closure properties of regular languages.
    
    \textit{Hint: Regular languages are closed under \emph{complement (i.e., $\Sigma^* \setminus C$)} and
    \emph{intersection ($A \cap B$)}. You may use these facts without proof.} 
\end{enumerate}

\textbf{Solution.} \textit{Write your solution here.}

\vspace{5em}

\section*{Problem 4: What the pumping lemma does not imply (20 points)}

Consider the language
\[
    F=\{a^i b^j c^k \mid i,j,k\geq 0,\ \text{and if } i=1 \text{ then } j=k \}.
\]
Namely, the numbers of $b$'s and $c$'s must be equal only when there is exactly one $a$.

\begin{enumerate}
    \item[(a)] \textbf{(10 points)}
    Show that $F$ satisfies the pumping property of regular languages:
    there exists a positive integer $p$ such that every $s \in F$ with $|s| \geq p$ has a decomposition $s=xyz$ satisfying
    \[
        |y|>0, \qquad |xy| \leq p, \qquad xy^t z \in F \quad \text{for every integer } t\geq 0.
    \]

    \item[(b)] \textbf{(8 points)} Prove that $F$ is not regular.
    You may use closure properties established in lecture or a direct argument about DFAs. 

    \item[(c)] \textbf{(2 points)} Explain why parts (a) and (b) do not contradict the pumping lemma.
\end{enumerate}

\textbf{Solution.} \textit{Write your solution here.}

\vspace{5em} 

\section*{Problem 5: A CFG and a PDA (20 points)}

Let $\Sigma=\{0,1,2\}$ and
\[
    D=\{0^i1^j2^k \mid i,j,k\geq 0,\ i=j\text{ or }j=k\}.
\]
Here ``or'' is inclusive: a string may satisfy either equality or both.

\begin{enumerate}
    \item[(a)] \textbf{(8 points)}
    Give a context-free grammar generating $D$.
    Explain why your grammar generates $D$.
    
    \item[(b)] \textbf{(8 points)}
    Give a PDA recognizing $D$.  Describe how the PDA works. 
    Give a diagram or a transition table. In either case, specify the stack alphabet, the nondeterministic choices, the stack operations, and when acceptance occurs. 
    Explain why your PDA recognizes $D$.

    \item[(c)] \textbf{(4 points)}
    Show that the context-free grammar you gave in part (a) is ambiguous by giving two distinct parse trees for one string in $D$. 
\end{enumerate}

\textbf{Solution.} \textit{Write your solution here.}

\vspace{5em}

\section*{6: Survey (not graded)}

Did you use AI for this homework assignment?
If yes, how did you use it?

(This question is not graded. The purpose of this question is to help you reflect on how AI can help your learning and research.)

\textbf{Response.} \textit{Write your response here.}

\end{document}
